%%%%%%%%%%%%
%
% $Autor: Wings $
% $Datum: 2019-03-05 08:03:15Z $
% $Pfad: Automatisierung/Skript/Produktspezifikation/Powerpoint/AMF.tex $
% $Version: 4250 $
% !TeX spellcheck = en_GB/de_DE
% !TeX encoding = utf8
% !TeX root = filename 
% !TeX TXS-program:bibliography = txs:///biber
%
%%%%%%%%%%%%

%Quelle: https://learn.sparkfun.com/tutorials/efficient-arduino-programming-with-arduino-cli-and-visual-studio-code/introduction-to-the-arduino-cli

%todo https://itp.nyu.edu/physcomp/resources/using-the-arduino-command-line-interface/


\chapter{Kommandozeilenschnittstelle CLI}

Die \ac{cli} ist eine Weiterentwicklung der Firma Arduino. Alle Funktionen, die über die Arduino IDE ausgeführt werden, können auch auf der Ebene einer Kommandozeile aufgerufen und durchgeführt werden. \cite{ArduinoCLI:2022}. Der Quellcode  kann frei genutzt werden, allerdings muss man bei einer gewerblichen Nutzung eine Lizenz bei der Firma anfragen \cite{ArduinoCLIGit:2022}. Eine ausführliche Dokumentation wird zur Verfügung gestellt. \cite{ArduinoCLIIntro:2020,ArduinoCLIDoc:2022}

Bei der Entwicklung der \ac{cli} stehen drei Aspekte im Fokus. Einerseits ermöglicht die Schnittstelle die Einbettung der Entwicklung in die gewohnte Entwicklungsumgebung. So kann die Schnittstelle auch mit  Edge Impulse\index{Edge Impulse} genutzt werden. Zweitens ermöglicht sie die Integration der Entwicklung in die Prozesse Continuous Development\index{Continuous Development} and Continuous Integration\index{Continuous Integration}. Sie ermöglicht damit die Automatisierung der typischen Aktivitäten im Bereich der Softwareentwicklung. Des Weiteren ermöglicht es nun die vereinfachte Kommunikation mit Edge Computern, zum Beispiel mit einem Raspberry PI.



\section{Installation der Kommandozeilenschnittstelle CLI}

Da \ac{cli} ständig weiterentwickelt wird, muss zunächst die aktuelle Version aus dem GitHub-Projekt \url{https://arduino.github.io/arduino-cli/0.22/installation#latest-release}, 
siehe \cite{ArduinoCLIGit:2022},  geladen werden.
%% https://github.com/arduino/arduino-cli#download-the-latest-unstable-alpha-preview.
Die heruntergeladene zip-Datei mit der Bezeichnung \FILE{arduino-cli\_0.22.1\_Windows\_64bit} in dem gewünschten Ordner. Dieser Pfad muss in der Liste der Umgebungsvariablen auf dem Computer hinzugefügt werden. Nachdem der Pfad hinzugefügt wurde, muss in einer Kommandozeile  das Kommando \SHELL{arduino-cli} aufgerufen werden. Die Installation startet dann. Nachdem die Installation abgeschlossen ist, kann \SHELL{arduino-cli board list} in der Eingabeaufforderung eingegeben werden. Das angeschlossene Gerät wird als Arduino Portenta H7\index{Portenta H7} mit Portnummer und Typ angezeigt, wie in Abbildung \ref{ArduinoInstallation} dargestellt. Es ist zu beachten, dass die Versionsnummer, hier 0.22.1, gegebenenfalls anzupassen ist. Wenn eine Nicla Vision\index{Nicla Vision} angeschlossen ist, werden entsprechende Meldungen ausgegeben.
Das Programm \FILE{arduino-cli-0.22.1.windows.exe} wird nun in \FILE{arduino-cli.exe} umbenannt. Daher kann im Folgenden der Programmname \FILE{arduino-cli.exe} verwendet werden.

\begin{figure}[H]
	\centering
	\includegraphics[width=\textwidth]{Arduino/arduino-cli}
	\caption{Installation von Arduino-CLI  }
	\label{ArduinoInstallation}
\end{figure}


\section{Konfiguration der Kommandozeilenschnittstelle CLI}


%todo https://learn.sparkfun.com/tutorials/efficient-arduino-programming-with-arduino-cli-and-visual-studio-code/introduction-to-the-arduino-cli

%Vielleicht ist es nur ein Alpha-Bug, aber das Arduino CLI auf meinem Rechner hat Schwierigkeiten, meine Arduino Sketchbook- und Boardmanager-Installationen zu finden. 
Damit \ac{cli} dir Arduino-Installation finden kann, hilft es, eine Konfigurationsdatei für \ac{cli} zu erstellen. Diese Konfigurationsdatei ist in einem YAML-Format definiert.
Zur Erstellung einer Basis-Konfigurationsdatei kann \ac{cli}  verwendet werden, in dem in einer Kommandozeile 

\medskip

\SHELL{arduino-cli.exe config init}

\medskip

eingegeben wird. Dieser Befehl erstellt eine neue Datei \FILE{.cli-config.yml}. Dort sind alle notwendigen Parameter deklariert. Folgende Parameter müssen in der Regel geändert werden:

\begin{itemize}
	\item \SHELL{sketchbook\_path}: Verzeichnis des Arduino-Sketche. Hier werden alle Bibliotheken und Hardware-Definitionen installiert.
	\item \FILE{arduino\_data}:  Installationsort des Arduino-Boards und des Bibliotheksmanagers. In den meisten Fällen sollte dies nicht geändert werden müssen.
\end{itemize}


Die anderen Optionen können in der Regel auf ihren Standardwerten bleiben.




\section{Verwendung der Kommandozeilenschnittstelle CLI}

Die Datei \FILE{README.md} im Arduino CLI GitHub Repository enthält eine hervorragende Übersicht über die Funktionen und Möglichkeiten.\cite{ArduinoCLIGit:2022} 
Es ist daher empfohlen, sie aus zu lesen, bevor man  fortfährt. 

\medskip

Hier sind einige wichtige Möglichkeiten:


\subsection{Erkennung der angeschlossenen Boards}




Mit Hilfe von \ac{cli} kann abgefragt werden, welche Boards installiert sind:

\medskip


\SHELL{arduino-cli board listall}

\medskip

Falls das angeschlossenes Board vermisst wird, muss es installiert werden:

\medskip

\SHELL{arduino-cli core install arduino:avr}

\medskip

Mit Hilfe von \ac{cli} kann abgefragt werden, welche Boards angeschlossen sind:

\medskip


\SHELL{arduino-cli board list}



\subsection{Erstellung eines neuen Sketches}

Der erste Schritt ist die Erstellung eines neuen Sketches:

\medskip

\SHELL{arduino-cli sketch new cli\_test}

\medskip

Dieser Befehl erstellt ein Verzeichnis mit dem Namen \FILE{cli\_test}, die eine Datei mit dem gleichen Namen enthält.

\subsection{Kompilieren eines Sketches}

Die Compile-Funktion von \ac{cli} kann verwendet werden, um einen Sketch für jedes unterstützte Board zu kompilieren. Die entscheidende Option, die diese Funktion benötigt, ist der Board-Typ, der mit der Option \SHELL{--fqbn} angegeben werden kann. Die Abkürzung \SHELL{fqbn} bedeutet \glqq fully-qualified board name\grqq, was \glqq vollqualifizierter Board-Name\grqq{} bedeutet. Beispielsweise stehen folgende Boards zur Verfügung:


\begin{itemize}
  \item Arduino Uno: \SHELL{arduino:avr:uno}
  \item Arduino Mega: \SHELL{arduino:avr:mega}
  \item SparkFun RedBoard oder SparkFun BlackBoard: \SHELL{SparkFun:avr:RedBoard}; dies erfordert die zusätzliche Installation der Definition \glqq SparkFun avr board\grqq.
  \item  SparkFun SAMD21 Mini: \SHELL{SparkFun:samd:samd21\_mini}; dies erfordert die zusätzliche Installation der Definition \glqq SparkFun samd board\grqq.
  \Mynote{Portenta H7 und Nicla Vision fehlen}
\end{itemize}

Die möglichen Boards sind wie folgt aufgebaut:  \SHELL{manufacturer:architecture:board}.

\bigskip

Mit dem folgenden Kommando kann der Beispiel-Sketch, der sich im Ordner  \PATH{C:/Users/user.name/Documents/Arduino/} befindet, für einen Arduino UNO kompiliert werden.:

\medskip

\SHELL{arduino-cli compile --fqbn arduino:avr:uno C:/Users/user.name/Documents/Arduino/cli\_test}

\medskip

In der Abbildung \ref{ArduinoCompile} sieht man die Meldungen von \ac{cli}.

\begin{figure}[H]
	\centering
	\includegraphics[width=\textwidth]{Arduino/arduino-cli-compile-cmd.png}
	\caption{Kompilieren eines Sketches mit \ac{cli}}
	\label{ArduinoCompile}
\end{figure}


Es können folgende Flags hinzugefügt werden:

\begin{itemize}
	\item Ausführlich \SHELL{-v}:  Nützlich, wenn alle Optionen und Dateien angezeigt werden sollen, die in Ihrem Sketch kompiliert werden.
	\item Build Pfad \SHELL{--build-path [string]}: Nützlich, wenn die kompilierten Objekt- und Hex-Dateien zu speichern sind. Auf meinem Windows-System muss der Wert dieses Parameters ein vollständiger Pfad sein.	
\end{itemize}

\subsection{Hochladen eines Sketches}

Ein kompilerter Sketch kann hochgeladen werden. Wie der Kompilierbefehl erfordert auch der Upload-Befehl eine \SHELL{--fqbn}. Außerdem wird eine serielle Schnittstelle zum Hochladen benötigt, die mit der Option \SHELL{-p} festgelegt wird.
Der folgende Befehl lädt den Beispiel-Sketch an einen Windows-COM-Anschluss auf COM18 hoch:

\medskip

\SHELL{arduino-cli upload -p COM18 --fqbn -v arduino:avr:uno C:/Users/user.name/Documents/Arduino/cli\_test}

\medskip


Falls es ordnungsgemäß durchgeführt wurde, sollte die RX/TX-LEDs des Arduinos zu blinken beginnen und kurz darauf einen leeren Sketch auszuführen.


\subsection{Installation von Bibliothelen}

Falls eine Bibliothek benötigt wird, so kann zunächst mit dem Kommando 

\medskip

\SHELL{arduino-cli lib search ethernet}

\medskip

on die Bibliothek, hier \SHELL{ethernet}, installiert ist. Mit dem Befehl

\medskip

\SHELL{arduino-cli lib install "{}UIPEthernet"}

\medskip

wird sie dann installiert.